\chapter*{Conclusion générale}
Ce stage a permis de structurer une vision cohérente de la transformation d'un WMS traditionnel vers un Smart WMS. Le travail ne s'est pas limité à identifier des technologies avancées ; il a consisté à replacer chaque capacité dans un besoin métier, une condition d'activation et une architecture fonctionnelle claire.

Le benchmark a montré trois principes complémentaires : l'unification de l'exécution, la richesse du modèle métier et l'orchestration des tâches. L'analyse des besoins a ensuite permis de formaliser les use cases Core, Smart et Optional. L'architecture fonctionnelle cible positionne les systèmes externes, le cœur opérationnel, les tâches, les données, les ressources, l'intelligence et les terminaux terrain dans une structure modulaire.

La suite du travail consiste à consolider la modélisation UML, valider les scénarios avec des acteurs métier et transformer progressivement l'architecture logique en prototype plus complet. L'apport principal du stage réside dans cette clarification : un Smart WMS robuste n'est pas un WMS auquel on ajoute de l'IA partout, mais un système capable d'utiliser les bonnes données, au bon endroit, pour améliorer les décisions et l'exécution réelle.
